\chapter{Manual de usuario}
\label{app:manual_usuario}

\ph{Este apéndice está dirigido al usuario final del sistema. A diferencia
del manual de instalación, no asume conocimientos técnicos y se centra en
\emph{cómo usar} el sistema una vez en marcha. Incluye capturas de pantalla
de las acciones principales cuando completes la memoria final.}

\section{Acceso al sistema}
\label{sec:acceso}

\ph{Indica la URL (o forma de acceso) del sistema y los pasos para
registrarse e iniciar sesión.}

\ph{\begin{itemize}
  \item \ph{URL de acceso: \url{https://example.org}}
  \item \ph{Registro de cuenta nueva.}
  \item \ph{Inicio de sesión.}
  \item \ph{Recuperación de contraseña (si aplica).}
\end{itemize}}

\ph{(Espacio reservado para una captura de la pantalla de inicio.)}

El sistema puede utilizarse mediante dos alternativas. En primer lugar, es
posible ejecutar la aplicación en el entorno local siguiendo los pasos
descritos en la Sección~\ref{sec:verificacion-instalacion} del manual de
instalación y despliegue.

Como alternativa, se puede acceder al despliegue oficial disponible públicamente
mediante la siguiente dirección:

\begin{itemize}
    \item \url{https://www.uvlhub.io/}
\end{itemize}

La generación de modelos no requiere que el usuario disponga de una cuenta ni
que realice un proceso de autenticación previo.

La pantalla inicial de UVLHub se muestra en la Figura~\ref{fig:inicio-uvlhub}.

\begin{figure}[H]
    \centering
    \makebox[\textwidth][c]{%
        \includegraphics[width=1.2\textwidth]{inicio_uvlhub.png}
    }
    \caption{Pantalla inicial de UVLHub. Elaboración propia}
    \label{fig:inicio-uvlhub}
\end{figure}

\section{Funcionalidad principal}
\label{sec:funcionalidad_principal}

\ph{Describe paso a paso cómo realizar la tarea principal del sistema.
Conviene incluir capturas o anotaciones que destaquen los botones y campos
relevantes.}

\ph{\begin{enumerate}
  \item \ph{Paso 1: dónde hacer clic, qué introducir.}
  \item \ph{Paso 2: qué pantalla aparece y qué hacer en ella.}
  \item \ph{Paso 3: cómo confirmar la acción.}
  \item \ph{Paso 4: cómo verificar el resultado.}
\end{enumerate}}

La funcionalidad principal del sistema consiste en la generación automática de
modelos de características mediante un asistente de configuración
(\textit{wizard}). La estructura completa del flujo de generación y las
pantallas correspondientes a cada etapa pueden consultarse en la
Sección~\ref{subsec:mapa_navegacion}.

El funcionamiento del asistente consiste en configurar los diferentes
parámetros de generación según las necesidades del usuario. En cada etapa del
flujo se pueden modificar las opciones disponibles y avanzar al siguiente paso
mediante el botón \textit{Next}, situado en la parte inferior de la pantalla.
En caso de querer modificar una configuración anterior, es posible volver a
los pasos previos utilizando el botón \textit{Previous}.

Una vez completada la configuración, en el sexto y último paso del asistente se generan
los modelos de características solicitados. Desde esta pantalla el usuario
dispone de diferentes opciones sobre los modelos generados:
\begin{itemize}
    \item Descargar el conjunto de modelos generado mediante el botón
    \textit{Download}.
    \item Visualizar un modelo directamente desde la aplicación mediante el
    botón con el icono del ojo.
    \item Descargar individualmente cada modelo mediante el botón verde
    disponible junto a cada resultado.
    \item Abrir un modelo mediante el botón morado con el icono del
    cohete en FlamapyIDE, un entorno externo que permite visualizar, editar y
    realizar tareas de análisis automático sobre modelos de características
    (\url{https://github.com/flamapy/flamapy-ide}).
\end{itemize}

\section{Mensajes de error y resolución de incidencias}
\label{sec:errores}

\ph{Lista de errores frecuentes que puede encontrar el usuario, con su
significado y la acción recomendada para resolverlos.}

\ph{\begin{itemize}
  \item \textbf{\ph{Mensaje 1}}: \ph{causa habitual y solución.}
  \item \textbf{\ph{Mensaje 2}}: \ph{causa habitual y solución.}
  \item \textbf{\ph{Mensaje 3}}: \ph{causa habitual y solución.}
\end{itemize}}

Durante el proceso de configuración y generación, el sistema muestra mensajes
de error descriptivos cuando no se cumplen las reglas de negocio definidas en la Sección~\ref{sec:reglas_negocio},
indicando al usuario la condición que debe modificarse para poder continuar.

Existe un caso particular durante la generación de modelos en el que puede
aparecer un error no asociado a una configuración incorrecta, sino a una
exigencia excesiva sobre los recursos disponibles. Este problema puede ocurrir,
por ejemplo, al solicitar un número elevado de modelos con un gran número de
características, restricciones, variables por restricción o atributos.

Para solucionar esta situación se recomienda revisar los parámetros de
generación configurados, especialmente:
\begin{itemize}
    \item Número de modelos a generar.
    \item Número de características.
    \item Número de restricciones.
    \item Número de variables por restricción.
    \item Número de atributos.
\end{itemize}

Un ejemplo del mensaje mostrado en esta situación se presenta en la
Figura~\ref{fig:error-generacion-modelos}.

\begin{figure}[H]
    \centering
    \includegraphics[width=1\textwidth]{error_generacion_modelos.png}
    \caption{Ejemplo de error durante la generación de modelos. Elaboración propia}
    \label{fig:error-generacion-modelos}
\end{figure}

\section{Contacto y soporte}
\label{sec:contacto}

\ph{Vía de contacto para incidencias o soporte (correo, formulario,
repositorio de issues), si la hubiese.}

Para cualquier duda relacionada con el funcionamiento del sistema o para
reportar errores no contemplados en este manual, se puede contactar con el
desarrollador mediante las siguientes direcciones de correo electrónico:

\begin{itemize}
    \item antjimort@alum.us.es
    \item ajimenez34@us.es
\end{itemize}